\section{A context-dependent fold from 245 to 244 states}
\label{sec:contextual-fold}

Let $M$ be the old-label Routing245 table, with initial state $95$ and
target-only halt state $254$.  Construct $M'$ by deleting state $122$ and
making exactly the following four changes to retained transition rows:
\begin{center}
\begin{tabular}{ccll}
\hline
State & Read & Original $(\text{write},\text{move},\text{target})$
                 & Replacement\\
\hline
$92$  & $1$ & $(1,L,24)$  & $(1,R,66)$\\
$104$ & $1$ & $(1,R,122)$ & $(1,R,92)$\\
$106$ & $1$ & $(1,R,122)$ & $(1,R,92)$\\
$167$ & $1$ & $(1,L,92)$  & $(1,L,151)$\\
\hline
\end{tabular}
\end{center}
The deleted rows are $(122,0)\mapsto(0,R,25)$ and
$(122,1)\mapsto(1,R,66)$.  No other row is changed.
Consequently $M'$ has exactly $244$ working states, $488$ binary rows, and
the same unique halt target.  The initial state remains $95$.

The reduction is not a claim that states $92$ and $122$ have the same
behavior on arbitrary configurations.  Their entry contexts provide the
missing information.  In particular, an entry into state $122$ always
leaves a $1$ immediately to the left of the head.  The replacement can
consult that bit and return to the same position after two extra steps.
The row at $(167,1)$ is rerouted to an already existing backward branch so
that its former use of $(92,1)$ remains intact.

\subsection{An inductive condition on the original computation}

For a configuration $(q,h,T)$, where $T:\mathbb Z\to\{0,1\}$ is its tape,
define the condition $I$ by the following implications:
\[
\begin{aligned}
 q\in\{68,122,167\}&\quad\Longrightarrow\quad T(h-1)=1,\\
 q=81\ \text{and}\ T(h-1)=1
     &\quad\Longrightarrow\quad T(h-2)=1.
\end{aligned}
\]
This condition holds for every configuration whose state is $95$,
irrespective of its head position and tape.

It is preserved by every original transition.  Here is the complete
incoming-edge inventory for its four constrained states:
\begin{center}
\begin{tabular}{cl}
\hline
Target & Incoming rows of $M$\\
\hline
$68$  & $(52,1)\mapsto(1,R,68)$;
         $(123,1)\mapsto(1,R,68)$\\
$81$  & $(68,1)\mapsto(1,R,81)$;
         $(123,0)\mapsto(0,R,81)$\\
$122$ & $(104,1)\mapsto(1,R,122)$;
         $(106,1)\mapsto(1,R,122)$\\
$167$ & $(145,1)\mapsto(1,R,167)$;
         $(148,1)\mapsto(1,R,167)$\\
\hline
\end{tabular}
\end{center}
For targets $68,122,167$, the required bit is precisely the $1$ just
written by the incoming transition.  An entry into $81$ from $123$ writes
$0$ immediately to the left, making the antecedent false.  An entry from
$68$ writes the first required $1$ and inherits the second from $I$ at
state $68$.  All other target states impose no condition.  This proves
inductiveness without assumptions about source instructions, register
values, spare blocks, or the size of a header gap.

\subsection{Seven finite identities in arbitrary tape contexts}

Use all states other than $92$ and $122$ as permissible checkpoint
states. At a paired checkpoint take the whole prescribed block below;
intermediate visits to permissible states need not themselves be paired.  Except for four
entry rows, every original transition out of a checkpoint is an identical
replacement transition and reaches another checkpoint or the common halt
target.  The four exceptional entry rows are
$(104,1),(106,1),(167,1),(81,1)$.

For an entry at head position $h$, the following finite calculations give
equal endpoint state, head, and entire tape in the two machines.  Cells
outside the indicated windows are never read or written by the
calculations.  All statements therefore hold with arbitrary left and
right tape contexts.

\paragraph{The two former entries into state $122$.}
Let $q\in\{104,106\}$ and suppose $T(h)=1$.  If $T(h+1)=1$, the two
computations reach state $66$ at $h+2$ in two steps each and preserve the
tape.  If $T(h+1)=0$, the original computation reaches state $25$ at $h+2$
in two steps.  The replacement takes four steps:
\[
 q \xrightarrow{1}92\xrightarrow{0}107
   \xrightarrow{1}125\xrightarrow{0}25.
\]
Its head positions after these steps are respectively
$h+1,h,h+1,h+2$.  The visited bits are $1,0,1,0$, and every write restores
the bit just read.  Thus the complete endpoint agrees.  The third step
uses the $1$ written at $h$ by the first step; it needs no additional
invariant premise.  Splitting the two entry states and the two following
bits gives four finite identities.

\paragraph{The repaired entry from $167$.}
If $q=167$ and $T(h)=1$, condition $I$ gives $T(h-1)=1$.  The old two-step
path is $167,92,24$, and the new two-step path is $167,151,24$.
The relevant rows of $92$ in $M$ and $151$ in $M'$ both read $1$, write
$1$, move left, and enter $24$.  Both calculations therefore preserve
the tape and finish at state $24$ with head $h-2$.

\paragraph{The entry from $81$ with a preceding zero.}
If $q=81$, $T(h)=1$, and $T(h-1)=0$, both machines use
$81,92,107$, preserve the tape, and finish at head $h-2$ in two steps.
The unchanged zero branch of state $92$ suffices.

\paragraph{The entry from $81$ with a preceding one.}
If $q=81$ and $T(h)=T(h-1)=1$, condition $I$ gives $T(h-2)=1$.
The old state sequence is
\[
 81,92,24,37,55,20,2,
\]
and the new state sequence is
\[
 81,92,66,20,2.
\]
The old calculation takes six transitions; the new calculation takes
four.  Both finish in state $2$ at $h-2$.  Both change exactly the cells
$h-1$ and $h$ from $1$ to $0$, leaving $T(h-2)=1$ and every other cell
unchanged.  This is the only finite identity that uses the two-cell
implication at state $81$.

\subsection{Global halting equivalence}

\paragraph{Theorem.}
For every binary tape $T$ and head position $h$, the computations of $M$
and $M'$ started in $(95,h,T)$ have the same halting outcome.

\paragraph{Proof.}
State $95$ is a checkpoint and satisfies $I$ vacuously.  Follow the
original computation from one selected checkpoint to the next, using
the prescribed complete block at each exceptional entry.  At an ordinary
entry both machines make the same one-step transition.  At one of the
four exceptional entries, the appropriate finite identity above produces
an identical checkpoint configuration in both machines.  The original
condition $I$ remains true because it is invariant under each original
literal step.  Induction thus gives identical paired checkpoint
configurations for as long as the computation continues.

Every block consumes between one and six original steps and between one
and four replacement steps.  Hence infinitely many blocks consume
unbounded time in both computations.  No exceptional block contains a
halt state, and all ordinary halting transitions agree.  A halt in either
machine therefore occurs exactly when the other machine halts.  This
proves both directions, including nonhalting preservation.\hfill$\square$

This theorem applies directly to the all-zero tape.  It transfers any
proved calibration of Routing245 to the new exact table without changing
the encoded source program or the logical theory.  In particular, from
the established Routing245 theorem one obtains
\[
 \mathrm{ZFC}\vdash
 \neg\operatorname{Halt}(M',0^{\mathbb Z})
 \ \longleftrightarrow\ \operatorname{Con}(\mathrm{ZFC}).
\]
The finite transition equalities and induction on numerical computation
prefixes give an arithmetically formalizable reduction; a newly checked
Lean or Rocq endpoint for $M'$ is a separate artifact and is not asserted
by the Python proof checker.

\subsection{Exact executable evidence}

The portable checker \texttt{prove\_contextual\_fold.py} checks the exact
four-row patch and state inventory, verifies the complete predecessor
lists, checks the local Boolean invariant, and replays the seven finite
identities while rejecting any read outside the represented window.
The invariant implication uses only offsets $-3,-2,-1,0,1$; exhaustive
Boolean evaluation of that finite formula is a universal verification,
independent of any bounded run of the machine.

The checker verifies $482$ common one-step rows and seven exceptional
window cases.  Its evidence is in \texttt{contextual\_fold\_proof.json}.
This certificate has a substantially smaller conceptual scope than a
complete register-program interpreter proof: it proves a direct
simulation between two fixed Turing machines.

An additional independent sparse/dense replay gives the same complete
Routing245 root frame after $407\,951$ steps of $M'$.  Its head is $-5$,
its header is at $176$, the physical register vector is
$(0,0,0,3,0,0,3,0,0)$, and its spare list contains $167$ values equal to
$1$.  There are exactly $351$ one-cells.  This finite initialization is
consistent with the global simulation theorem, although the theorem
does not depend on the initialization calculation.
